\section*{अध्याय \ref{ch:b3:ocean-dynamics} --- महासागर गतिकी}

\begin{solution}{exo:b3:ocean-dynamics:1}
$f = \qty{9.37e-5}{s^{-1}}$; $\tau/(\rho f) = 0.12/(1025\times\num{9.37e-5})
= \qty{1.25}{m^2/s}$, पवन की दाईं ओर, दक्षिण की ओर
\qty{1000}{km} के पार, \qty{1.25e6}{m^3/s}, \qty{1.25}{\sverdrup}।
\end{solution}

\begin{solution}{exo:b3:ocean-dynamics:2}
$f = \qty{1.031e-4}{s^{-1}}$; $D_E = \pi\sqrt{2\times0.05/\num{1.031e-4}}
= \qty{98}{m}$। वहाँ की धारा सतही धारा की $\mathrm e^{-\pi} = \num{0.043}$ गुनी होती है,
और विपरीत दिशा में बहती है।
\end{solution}

\begin{solution}{exo:b3:ocean-dynamics:3}
$\beta = 2\Omega\cos\ang{30}/a = \qty{1.98e-11}{m^{-1}.s^{-1}}$; $V =
\num{-1e-7}/(1025\times\num{1.98e-11}) = \qty{-4.92}{m^2/s}$; \qty{5000}{km} के पार
\qty{-24.6}{\sverdrup}: दक्षिण की ओर।
\end{solution}

\begin{solution}{exo:b3:ocean-dynamics:4}
दोनों \omterm{def:b3:ocean-dynamics:tides}{ज्वारीय घटक} $1/(1/12 - 1/12.4206) =
\qty{354}{h}$ के बाद फिर से समान कला में आ जाते हैं: \num{14.77}\text{ दिन}, अर्थात् एक सिनोडिक मास का आधा।
\end{solution}

\begin{solution}{exo:b3:ocean-dynamics:5}
\ang{40} N पर, $f = \qty{9.37e-5}{s^{-1}}$, $\beta =
\qty{1.75e-11}{m^{-1}.s^{-1}}$, $\tau_x = \qty{0.05}{Pa}$ और
$\partial\tau_x/\partial y = 0.1/\num{1.11e6} = \qty{9.0e-8}{Pa/m}$ मिलता है।
पहला पद, $-\num{9.0e-8}/(1025\times\num{9.37e-5}) =
\qty{-9.38e-7}{m/s}$; दूसरा पद, $0.05\times\num{1.75e-11}/(1025\times
(\num{9.37e-5})^2) = \qty{+9.7e-8}{m/s}$. $w_E = \qty{-8.4e-7}{m/s}$,
\qty{-26.5}{m/yr}: पानी नीचे की ओर धकेला जाता है।
\end{solution}

\begin{solution}{exo:b3:ocean-dynamics:6}
$\delta_S = \num{2e-6}/\num{2e-11} = \qty{100}{km}$; गति $=
\num{30e6}/(\num{1e5}\times800) = \qty{0.375}{m/s}$।
\end{solution}

\begin{solution}{exo:b3:ocean-dynamics:7}
$\beta = \qty{1.98e-11}{m^{-1}.s^{-1}}$; $\delta_M = (10^4/\num{1.98e-11})^{1/3}
= \qty{79.6}{km}$, और जेट की चौड़ाई $3.628\,\delta_M = \qty{289}{km}$ है।
$A_H$ के आठ गुना कम होने पर दोनों आधे हो जाते हैं, क्योंकि $8^{1/3} = 2$।
\end{solution}

\begin{solution}{exo:b3:ocean-dynamics:8}
$f = \qty{8.57e-5}{s^{-1}}$; $v = (21.3 - 14.5)/(\num{8.57e-5}\times\num{8e4})
= \qty{0.99}{m/s}$, $B$ दाईं ओर होने से धनात्मक, $A$ से $B$ की ओर देखने पर
पूर्व की ओर, धारा उत्तर की ओर जाती है। समुद्र की सतह
$6.8/9.81 = \qty{0.69}{m}$ ऊँची रहती है, $B$ पर;
$B$ पर जल-स्तंभ हल्का और गर्म है, जो धारा की दाईं ओर है।
\end{solution}

\begin{solution}{exo:b3:ocean-dynamics:9}
वसंत-ज्वार परास $1.5K(1 + 0.46) = \qty{0.78}{m}$; लघु-ज्वार परास $1.5K(1 - 0.46)
= \qty{0.29}{m}$; अनुपात $1.46/0.54 = 2.7$।
\end{solution}

\begin{solution}{exo:b3:ocean-dynamics:10}
\cref{thm:b3:ocean-dynamics:kelvin} के अनुसार, $u = \pm(g/c)\eta$ के लिए $F(x \mp
ct)$, और $fu = -g\,\partial\eta/\partial y$ हो जाता है $\beta_0y(\pm g/c)G =
-gG'$, अतः $G = \exp(\mp\beta_0y^2/2c)$। केवल ऊपरी चिह्न, यानी पूर्व की ओर जाने वाली
तरंग भूमध्य रेखा के दोनों ओर सीमित रहती है। $\beta_0 =
\qty{2.29e-11}{m^{-1}.s^{-1}}$, $\sqrt{2.5/\num{2.29e-11}} = \qty{330}{km}$।
\end{solution}

\begin{solution}{exo:b3:ocean-dynamics:11}
(क)~तल से $\beta v = f\,\partial w/\partial z$ का समाकलन, जहाँ $w = 0$ है,
गहन परत के ऊपरी सिरे तक, जहाँ $w = W_0$ है: $\beta V = fW_0$, और $f/\beta =
a\tan\varphi$। (ख)~$V = \num{6.371e6}\times0.577\times10^{-7} =
\qty{0.368}{m^2/s}$, \qty{5000}{km} के पार \qty{1.84}{\sverdrup}, ध्रुव की ओर।
(ग)~यह उच्च-अक्षांश स्रोतों द्वारा पोषित गहन पश्चिमी सीमा धारा से आता है;
आंतरिक भाग स्रोतों की ओर बहता है, और सीमा धारा उनसे दूर की ओर
बहती है।
\end{solution}

\begin{solution}{exo:b3:ocean-dynamics:12}
(क)~तल से ऊपर की ओर $z$ मापने पर, $A_zW'' = \mathrm if(W - W_g)$,
दाब-प्रवणता वही है जो ऊपर के भू-विक्षेपी प्रवाह की है। $W - W_g =
C\mathrm e^{-(1+\mathrm i)z/\delta}$ ऊपर की ओर सीमित है, और $W(0) = 0$ से
$C = -W_g$ प्राप्त होता है। (ख)~परिवहन त्रुटि $\int_0^\infty(W - W_g)\,\mathrm dz
= -W_g\delta/(1+\mathrm i) = -\tfrac12\delta(1 - \mathrm i)W_g$, अर्थात्
$-\tfrac12\delta(u_g + v_g,\ v_g - u_g)$। भू-विक्षेपी प्रवाह
अपसारी-रहित होने से, $w = -\nabla\cdot(\text{त्रुटि}) = \tfrac12\delta(\partial_x
v_g - \partial_y u_g) = \tfrac12\delta\zeta_g$। (ग)~$\delta =
\sqrt{0.02/10^{-4}} = \qty{14.1}{m}$; $r = 10^{-4}\times14.1/8000 =
\qty{1.77e-7}{s^{-1}}$, अर्थात् \num{65}\text{ दिन} का स्पिन-डाउन समय।
\end{solution}

\begin{solution}{pb:b3:ocean-dynamics:1}
\textbf{1.} $\tau_x = \qty{-0.1}{Pa}$ (पुरवा), 0 और \qty{+0.1}{Pa}
(पछुआ)।

\textbf{2.} $f = \num{4.99e-5}$, \num{8.37e-5} और
\qty{1.117e-4}{s^{-1}}।

\textbf{3.} $V_E = -\tau_x/(\rho f)$: \qty{+1.96}{m^2/s} (उत्तर), \ang{20} N पर
और \qty{-0.87}{m^2/s} (दक्षिण), \ang{50} N पर। ये अभिसरित होते हैं।

\textbf{4.} $D_E = \pi\sqrt{2\times0.05/\num{8.37e-5}} = \qty{109}{m}$।

\textbf{5.} $\tau_x = 0$ \ang{35} N पर, अतः $w_E = -(1/\rho f)\,
\partial\tau_x/\partial y = -\pi\tau_0/(\rho fL) = \qty{-1.10e-6}{m/s}$,
\qty{-34.7}{m/yr}, नीचे की ओर।

\textbf{6.} $500/34.7 = \num{14.4}\text{ वर्ष}$।

\textbf{7.} $L = a\times\pi/6 = \qty{3336}{km}$; $\beta =
\qty{1.875e-11}{m^{-1}.s^{-1}}$।

\textbf{8.} $\operatorname{curl}_z\vec\tau = -\pi\tau_0/L =
\qty{-9.42e-8}{Pa/m}$।

\textbf{9.} $V = \num{-9.42e-8}/(1025\times\num{1.875e-11}) =
\qty{-4.90}{m^2/s}$, दक्षिण की ओर।

\textbf{10.} $4.90\times\num{6e6} = \qty{2.94e7}{m^3/s}$,
\qty{29.4}{\sverdrup} दक्षिण की ओर।

\textbf{11.} $\qty{4.9}{mm/s}$; \qty{1000}{km} \qty{2.0e8}{s} में,
\num{6.5}~वर्ष।

\textbf{12.} $fw_E/\beta = \num{8.37e-5}\times(\num{-1.10e-6})/
\num{1.875e-11} = \qty{-4.90}{m^2/s}$: वही, क्योंकि एकमन
परिवहन वहाँ शून्य होता है जहाँ $\tau_x = 0$ है। $w_E$ की दर से दबा हुआ एक स्तंभ
अपनी \omterm{def:b3:atmospheric-dynamics:pv}{विभव भ्रमिलता} $f/h$ केवल छोटे $f$ की ओर, यानी दक्षिण की ओर
खिसककर बनाए रखता है, $\beta V_g = fw_E$ की दर से।

\textbf{13.} केवल वही सीमा-परत संभव है जिसमें आंतरिक क्षेत्र को किया गया सुधार
दीवार से दूर की ओर क्षीण होता है; $\beta > 0$ के साथ यह
$\mathrm e^{-x/\delta_S}$ है, जो पूर्व की ओर क्षीण होता है, और यह
केवल पश्चिमी दीवार से ही होता है। भ्रमिलता के संदर्भ में, उत्तर की ओर बहती धारा में घर्षण
पूर्वी दीवार के विरुद्ध धनात्मक भ्रमिलता हटा देता, जो गलत चिह्न होता।

\textbf{14.} $\delta_S = \num{e-6}/\num{1.875e-11} = \qty{53.3}{km}$।

\textbf{15.} $A = \qty{4.90}{m^2/s}$, $X/\delta_S = 112.5$:
$4.90\times111.5 = \qty{546}{m^2/s}$, \qty{0.55}{m/s} का प्रवाह
\qty{1000}{m} के आर-पार।

\textbf{16.} जहाँ $(X/\delta_S)\mathrm e^{-x/\delta_S} = 1$: $x =
\delta_S\ln(X/\delta_S) = \qty{252}{km}$।

\textbf{17.} $A(X - \delta_S - x_0) = 4.90\times(6000 - 53 - 252)\times
10^3 = \qty{27.9}{\sverdrup}$। आंतरिक प्रवाह अब भी दक्षिण की ओर बह रहा है,
दीवार से सटी पट्टी में; कमी, $A(x_0 + \delta_S) =
\qty{1.5}{\sverdrup}$, वह आंतरिक परिवहन है जो
उसी पट्टी को पार कर गया होता।

\textbf{18.} $\delta_M = (\num{5e3}/\num{1.875e-11})^{1/3} = \qty{64.4}{km}$;
जेट $3.628\times64.4 = \qty{233}{km}$।

\textbf{19.} $1.163\times29.4 - A\times\qty{233}{km} = 34.2 - 1.1 =
\qty{33.0}{\sverdrup}$; अधिक भाग, जो आंतरिक परिवहन का \qty{16}{\percent} है,
जेट के ठीक पूर्व की प्रति-धारा में दक्षिण की ओर लौटता है।

\textbf{20.} $c = \sqrt{9.81\times4000} = \qty{198}{m/s}$; $L_R =
198/\num{8.37e-5} = \qty{2370}{km}$।

\textbf{21.} $198\times12.4206\times3600 = \qty{8860}{km}$।

\textbf{22.} परिमाप $2(6000 + 3336) = \qty{18670}{km}$, \qty{26.2}{h},
2.1 M$_2$ आवर्त।

\textbf{23.} वृत्त $2\pi a\cos\ang{35} = \qty{32790}{km}$ है, जिसे
एक चंद्र दिवस, \qty{24.84}{h} में तय करना है: \qty{367}{m/s}, जिसके लिए $c^2/g =
\qty{13.7}{km}$ जल-स्तंभ आवश्यक होता है।

\textbf{24.} जब रॉस्बी त्रिज्या बेसिन की चौड़ाई के तुल्य हो और
M$_2$ तरंगदैर्घ्य उसकी लंबाई के तुल्य हो, तो लगभग एक \omterm{def:b3:ocean-dynamics:amphidrome}{एम्फ़िड्रोमिक
बिंदु} बनता है, जो बेसिन के केंद्र के समीप होता है; उच्च जल उसके चारों ओर वामावर्त घूमता है,
तट को अपनी दाहिनी ओर रखते हुए।

\textbf{25.} $\tfrac34\times0.357\times\cos^2\ang{35} = \qty{0.18}{m}$।

\textbf{26.} पश्चिमी सीमा धारा उत्तर की ओर लौटाती है
\textbf{\qty[text-series-to-math=true]{29}{\sverdrup}}, जिसे आंतरिक क्षेत्र
दक्षिण की ओर ले जाता है (गणना के अनुसार \qty{29.4}{\sverdrup})।
\end{solution}
